\documentclass{article}

\textwidth=6.5in
\textheight=8.5in
\voffset=-54pt
\oddsidemargin=0in
\evensidemargin=0in
\setlength{\parskip}{8pt}
\setlength{\parindent}{0pt}

\usepackage{amsmath, amssymb}
\usepackage{ifthen}

\usepackage[inline]{enumitem}
\setlist{topsep=0pt, parsep=8pt}

\usepackage[rgb,table]{xcolor}\convertcolorsUfalse
\usepackage{graphicx}

% change hyperref options
\PassOptionsToPackage{hyphens}{url}
\usepackage[bookmarks,colorlinks,linkcolor=black,citecolor=blue,pdfstartview=FitH,urlcolor=blue]{hyperref}
\hypersetup{pdfpagemode=UseNone}
\usepackage{bookmark}

\usepackage{graphicx}
\usepackage{multicol}
\usepackage{framed}
\usepackage{array}

\usepackage{icomma}

\usepackage{marginnote}
\usepackage[colorinlistoftodos]{todonotes}
\renewcommand{\marginpar}{\marginnote}

\usepackage{soul}

\usepackage{pgfplots}
\pgfplotsset{compat=1.18}  
\pgfplotscreateplotcyclelist{mycolor}{black, blue, red, green!70!black, magenta, purple, cyan, orange }  
\pgfplotsset{
  disabledatascaling,
  cycle list=tikzcolor, 
  axis lines=middle,
  every axis plot/.style={no marks, thick},
  every axis/.style={
    enlarge x limits=false,
    enlarge y limits=auto,
    xlabel style={at={(current axis.right of origin)},anchor=west},
    ylabel style={at={(current axis.above origin)},anchor=south},
    % extra x and y ticks for asymptotes
    extra x tick style={grid=major, grid style={dashed,black}},
    extra y tick style={grid=major, grid style={dashed,black}},
    /pgf/number format/1000 sep={},
  },    
  tick label style = {font=\footnotesize, fill=\currentbgcolor, inner sep=0pt},    
  xticklabel shift=1pt,    
  scaled ticks=false,
  scaled y ticks=false,
  unbounded coords=jump,
  samples=101,
  minor tick length=0,
  scatter/classes={
    open={mark=*,draw=black,fill=white, mark size=2, thin},
    closed={mark=*,draw=black,fill=black, mark size=2, thin}
  },
  width=3in,
}
\pgfplotsset{open/.style={only marks, mark=*, fill=white, thin, mark size=2}}
\pgfplotsset{closed/.style={only marks, mark=*, draw=black, fill=black,
    thin, mark size=2}}
\pgfplotsset{labeled/.style={nodes near coords, only marks, mark=*, draw=black, fill=black, thin, mark size=2, point meta=explicit symbolic}}

\usepackage{tikz}
\usetikzlibrary{
  matrix,%
  % enables the snake paths
  decorations.pathmorphing,%
  % enables bigger arrows
  decorations.markings,%
  decorations.pathreplacing,%
  decorations.text,%
  calc,%
  patterns,%
  shapes.geometric,%
  arrows,
  decorations.shapes,
  positioning,
  plotmarks,
  intersections
}
\tikzset{>=latex} % sets default arrow type in tikz
\tikzset{font=\normalsize}
\tikzset{mark size=1.5pt, mark options=thin}
\tikzset{pin distance=4pt,
  every pin edge/.style={<-, thin, shorten <= -2pt}}

\interfootnotelinepenalty=10000
\renewcommand{\thefootnote}{(\arabic{footnote})}

\usepackage{multirow, bigdelim}

% change to \usepackage[sols]{handout} to see solutions
\usepackage[sols, grading, workshop]{handout}
\renewcommand{\workshopnote}{CA Notes.}    

\newcommand\iscurrentenumi[1]{%
  \ifthenelse{\equal{#1}{\theenumi}}{}{#1}}
\setenumerate[1]{ref=\#\arabic*}
\setenumerate[2]{ref=\protect\iscurrentenumi{\theenumi}(\alph*)}
\newcommand\enumiiprefix[2]{%
  \ifthenelse{{\equal{#1}{\theenumi}}\AND{\equal{#2}{(\alph{enumii})}}}{}{}%
  \ifthenelse{{\equal{#1}{\theenumi}}\AND{\NOT{\equal{#2}{(\alph{enumii})}}}}{#2}{}%
  \ifthenelse{\equal{#1}{\theenumi}}{}{#1#2}%
}
\setenumerate[3]{ref=\protect\enumiiprefix{\theenumi}{(\alph{enumii})}\roman*}

\makeatletter

% underline links               
\let\@href=\href
\renewcommand{\href}[2]{\@href{#1}{\ul{#2}}}

\makeatother

\definecolor{purple}{rgb}{0.5,0,1.0}

\newcounter{imagecounter}
\newcounter{columncounter}
\newenvironment{imagetable}[3][]{%
  \setcounter{imagecounter}{0}
  \setcounter{columncounter}{#2}
  \addtocounter{columncounter}{-1}
  \newcolumntype{i}{>{\raisebox{-1ex-\height}{\makebox[1em]{\hfill\stepcounter{imagecounter}(#3{imagecounter})}}\hspace{4pt}\begin{minipage}[t]{\linewidth/#2-1em-\tabcolsep*3}\arraybackslash\vspace{0pt}}l<{\end{minipage}}}
\begin{tabular}{i*{\value{columncounter}}{#1i}}}
{\end{tabular}}  

\newcommand{\doedfinity}[2][\newpage]{
  \begin{problem}[#1]
    Do the Edfinity assignment ``Problem Set #2''.

    We encourage you to write down any scratch work here, as well as
    things you want your future self to remember. Nothing you write
    here will be graded; it's just for you!
  \end{problem}
}

\newcommand{\psrnewpage}{\newpage}
\newcommand{\psreflection}[2][]{

  \ifthenelse{\not\equal{#1}{nocite}}{
    \begin{enumerate}[resume]
    \item
      \begin{problem}[1in]
        Please cite any help you got on this problem set:
        \begin{itemize}[nosep]
        \item If you worked with other students, please list their names here.
        \item If you used resources other than those on our Canvas site, please list them here.
        \end{itemize}
      \end{problem}

    \end{enumerate}
  }

  \probonly{%
    #2

    \psrnewpage

    \emph{Reflection space.} It's a good habit to take a few minutes at the end of each pset to reflect on what you learned and jot down notes to your future self. Here are a few reflection questions to get you started. Nothing you write here will be graded; it's just for you!
    \begin{itemize}
    \item Take a few minutes to look back at the problems. How could you explain to someone else how to do each problem? What was your strategy, and how did you know to use that strategy? If there are problems where you don't feel able to answer this, write down which ones they are. Come back to them in a few days when we post the homework solutions!

      \vspace{1.5in}

    \item Take a look back at the learning objectives on today's worksheet; how are you feeling about each one? If there are ones you know you need more practice with, write those down here.

      \vspace{1.5in}

    \item What did you find most challenging about this material? Do you have any lingering questions about it? If so, write them down here so you don't forget them. At some point, stop by office hours or MQC to ask!

      \vfill
    \end{itemize}      
  }
}

\usepackage{fancyhdr}
\lhead{\textcolor{white}{This content is protected and may not be shared, uploaded, or distributed.}}
\rhead{}
\rfoot{\vspace*{\fill}\footnotesize\textcopyright{} \emph{Harvard Math Ma}}
\renewcommand{\headrulewidth}{0pt}
\pagestyle{fancy}
\begin{document}

\begin{center}
  \large\textsc{Problem Set 35 - Cubics and Higher-Degree Polynomials}
\end{center}

\probonly{\emph{Note:} This is the second-to-last problem set!}

\begin{enumerate}
\item  \note{
    \begin{itemize}[nosep]
    \item Decide whether functions (given by formulas) are polynomials; if so, identify degree.
    \item Match quadratics to graphs
    \item Match factored polynomials to graphs
    \item Given pictures of functions, decide which are polynomials of odd degree and which are polynomials of even degree (includes non-polynomials that can be distinguished because of end behavior).
    \item Some true/false statements. 
    \end{itemize}
  }

  \doedfinity{35}

\item  

  \note{This is to show students different possible shapes of cubics.}

  \begin{enumerate}
  \item Two of the following three functions are cubics. For each cubic, find the number of critical points and the number of turning points. Explain your reasoning.

    For the function that isn't a cubic, simply write ``not a cubic''.

    \begin{enumerate}
    \item

      \begin{grading}
        1.5 points: 0.5 for calculating $a'$, 0.5 for recognizing there are no critical points, 0.5 for knowing there are therefore no turning points
      \end{grading}

      \begin{problem}[\vfill]
        $a(x) = x (x^2 + 7)$
      \end{problem}

      \begin{solution}
        This is a cubic. $a(x)=x^3+7x$, so $a'(x)=3x^2+7$. Since $a'(x)$
        has no real roots, $a(x)$ has no critical points and no turning
        points.
      \end{solution}

    \item

      \begin{grading}
        2.5 points:
        \begin{itemize}[itemsep=0pt, topsep=0pt]
        \item 1 for finding the critical point
        \item 1.5 for identifying that it isn't a turning point.

          The Second Derivative Test is inconclusive, so if they use that to justify this, deduct 1.5. They can use either the First Derivative Test or the fact that a cubic must have an even number of turning points (because of the end behavior).
        \end{itemize}
      \end{grading}

      \begin{problem}[\vfill]
        $b(x) = -3x + 6x^2 - 4x^3$
      \end{problem}

      \begin{solution}
        This is a cubic. $b'(x)=-3+12x-12x^2$. We can factor
        this as
        \begin{align*}
          b'(x) &= -3(4x^2 -4x + 1) \\
          &= -3(2x-1)^2
        \end{align*}
        so $b(x)$ has a critical point at $x=\frac{1}{2}$. It
        is not a turning point because $b'(x)$ does not change sign
        at $x=\frac{1}{2}$ (also because a cubic cannot have an odd
        number of turning points.)
      \end{solution}

    \item

      \begin{grading}
        0.5 points
      \end{grading}

      \begin{problem}[\nstretch{0.75}]
        $c(x) = x^3 + 7 x^{-1}$
      \end{problem}

      \begin{solution}
        This is not a cubic.
      \end{solution}

    \end{enumerate}

  \item

    \begin{grading}
      2 points, 1 per cubic. A possible breakdown for each cubic is: 0.5 for identifying which graphs match what they found in (a) (there should be 2), and 0.5 for justifying which of those 2 is the right one (end behavior is the most obvious, but they could also do something like test a point).
    \end{grading}

    \begin{problem}
      For each of the two cubics in (a), pick the picture below that shows its graph. Explain your choices briefly. 
    \end{problem}

    \begin{solution}
      \begin{itemize}
        \item $a(x)$ has no critical points, so the only possibilities
        are (C) and (F). Since the leading coefficient of $a(x)$ is
        positive, that leaves \fbox{(C)}.
        \item $b(x)$ has exactly one critical point, so
        the only possibilities are (B) and (E). Since the leading
        coefficient of $b(x)$ is negative, that leaves \fbox{(E)}.
      \end{itemize}
    \end{solution}

    \begin{imagetable}{3}{\Alph}
      \begin{tikzpicture}
        \begin{axis}[width=1.75in, ytick=\empty, ymax=10]
          \addplot+[domain=-3:3] {x^3 + 2*x^2}; 
        \end{axis}
      \end{tikzpicture} &
      \begin{tikzpicture}
        \begin{axis}[width=1.75in, ytick=\empty, ymin=-10]
          \addplot+[domain=-2:2] {-(-3*x + 6*x^2 - 4*x^3)}; 
        \end{axis}
      \end{tikzpicture} &
      \begin{tikzpicture}
        \begin{axis}[width=1.75in, ytick=\empty]
          \addplot+[domain=-3:3] {x*(x^2 + 7)}; 
        \end{axis}
      \end{tikzpicture}
      \vspace{18pt} \\
      \begin{tikzpicture}
        \begin{axis}[width=1.75in, ytick=\empty, ymin=-10]
          \addplot+[domain=-3:3] {-(x^3 + 2*x^2)}; 
        \end{axis}
      \end{tikzpicture} &
      \begin{tikzpicture}
        \begin{axis}[width=1.75in, ytick=\empty, ymax=10]
          \addplot+[domain=-2:2] {-3*x + 6*x^2 - 4*x^3}; 
        \end{axis}
      \end{tikzpicture} &
      \begin{tikzpicture}
        \begin{axis}[width=1.75in, ytick=\empty]
          \addplot+[domain=-3:3] {-x*(x^2 + 7)}; 
        \end{axis}
      \end{tikzpicture} 
    \end{imagetable}

    \pspace{\vfill}

  \end{enumerate}

\item 

  \begin{grading}
    2 points: 1 for explaining why we can eliminate A and C, 1 for explaining why we can eliminate D.
  \end{grading}

  \begin{problem}[\vfill\newpage]
    One of the following pictures shows the graph of a degree 4 polynomial. Which one? Explain how you eliminated the other choices.

    \begin{imagetable}{4}{\Alph}
      \begin{tikzpicture}
        \begin{axis}[xtick=\empty, ytick=\empty, enlarge y limits=false, width=1.5in]
          \addplot+[domain=-2:1.1] {(x+1)^2*(x^3+x)-1};
        \end{axis}
      \end{tikzpicture}
      &
      \begin{tikzpicture}
        \begin{axis}[xtick=\empty, ytick=\empty, enlarge y limits=false, width=1.5in, ymax=1]
          \addplot+[domain=-1.3:3] {-(3/4 - x + 2*x^2 - x^3 + x^4/4)}; 
        \end{axis}
      \end{tikzpicture}
      &
      \begin{tikzpicture}
        \begin{axis}[xtick=\empty, ytick=\empty, enlarge y limits=false, width=1.5in]
          \addplot+[domain=-2.7:2.4] {1.69792 - (-2 + x)*(-1 + x)*x*(1 + x)*(2.5 + x)}; 
        \end{axis}
      \end{tikzpicture}
      &
      \begin{tikzpicture}
        \begin{axis}[xtick=\empty, ytick=\empty, enlarge y limits=false, width=1.5in]
          \addplot+[domain=-2.2:2.6, restrict y to domain=-20:8] {-5 + (-2 + x)^2*(-1 + x)*x*(1 + x)*(2 + x)}; 
        \end{axis}
      \end{tikzpicture}
    \end{imagetable} 

  \end{problem}

  \begin{solution}
    The degree is even, so the ends must go in the same direction.
    This eliminates choices (A) and (C). A degree 4 polynomial
    can have at most 3 turning points, which eliminates choice (D),
    leaving only choice \fbox{(B)}.
  \end{solution}

\item  

  \note{The point of this problem is just to show explicitly that, for a cubic, the turning points aren't half-way between the zeros.}

  \begin{enumerate}
  \item

    \begin{grading}
      2 points:
      \begin{itemize}[nosep]
      \item 0.5 for having something of the form $a (x + 3)x (x - 5)$
      \item 0.5 for having a formula that satisfies $\displaystyle \lim_{x \to \infty} f(x) = -\infty$
      \item 1 for explanation
      \end{itemize}
    \end{grading}

    \begin{problem}[\stretch{0.35}]
      Write a formula for a cubic that has roots at $x = -3, 0, 5$ and has $\displaystyle \lim_{x \to \infty} f(x) = - \infty$. Explain how you came up with your formula. 
    \end{problem}

    \begin{solution}
      Since the roots are $x=-3, 0, 5$, the formula has to be of the
      form $ax(x+3)(x-5)$. To satisfy 
      $\lim\limits_{x \to \infty} f(x) = - \infty$, $a$ has to be negative,
      so one possible formula is \fbox{$-x(x+3)(x-5)$}.
    \end{solution}

  \item

    \begin{grading}
      4 points: 2 for finding critical points, 2 for classifying them and explaining clearly how they did so
    \end{grading}

    \begin{problem}[\vfill]
      Let $f(x)$ be your answer to (a). Find the critical points of $f$, and classify each as a local minimum, local maximum, or neither. Make sure to clearly communicate your justification; use words to let the reader know how you're classifying each critical point.
    \end{problem}

    \begin{solution}
      $f(x)=-x(x+3)(x-5)=-x^3+2x^2+15x$, so $f'(x)=-3x^2+4x+15$. This factors as
      \begin{align*}
        f'(x) &= -(3x^2-4x-15) \\
        &= -(3x+5)(x-3)
      \end{align*}
      so the critical points are $x=-5/3$ and $x=3$. Since $f'(x)$ is
      a quadratic with a negative leading coefficient, $f'(x)$ is negative
      for $x<-5/3$ and $x>3$ and positive for $-5/3<x<3$. This means
      that $x=-5/3$ is a local min and $x=3$ is a local max.
    \end{solution}

  \item

    \begin{grading}
      1 point
    \end{grading}

    \begin{problem}[\stretch{0.25}]
      Are the turning points of $f$ half-way between the roots of $f$? 
    \end{problem}

    \begin{solution}
      No. $x=-5/3$ is not half-way between $x=-3$ and $x=0$, and
      $x=3$ is not half-way between $x=0$ and $x=5$.
    \end{solution}

  \item

    \begin{grading}
      1 point
    \end{grading}

    \begin{problem}[\nstretch{0.35}]
      Does $f$ have a global maximum on $(-\infty, \infty)$? How about a global minimum? Explain how you know. 
    \end{problem}

    \begin{solution}
      $f$ has neither a global maximum nor a global minimum on $(-\infty,
      \infty)$. This is because $\lim\limits_{x\to\infty}f(x)=-\infty$ and 
      $\lim\limits_{x\to-\infty} f(x)=\infty$.
    \end{solution}

  \end{enumerate}

\item In each part, give an example of a polynomial $p(x)$ with the specified characteristics. Your answer should be a formula. Make sure to explain how you came up with your answer; the heart of each is figuring out a simple plan for finding such a polynomial. 

  \begin{workshop}
    The goal of these problem is for students to come up with a strategy, so don't just tell them how to approach these. A good question to ask them is: ``What are other ways you could say the given information?'' (For example, saying $p(x)$ has critical points at $x = 0, \pm 2$ is the same as saying $p'(x) = 0$ at those points.) Or: ``Can you say anything about $p'$ or $p''$?''
  \end{workshop}

  \begin{enumerate}
  \item

    \begin{grading}
      2 points:
      \begin{itemize}[nosep]
      \item 0.5 for giving an answer with the correct roots
      \item 1 for giving an answer with even degree and explaining why that's necessary
      \item 0.5 for an answer with $\displaystyle \lim_{x \to \infty} p(x) = -\infty$
      \end{itemize}
    \end{grading}

    \begin{problem}[\vfill]
      The only roots of $p(x)$ are $x = 2$, $x = 5$, and $x = 7$, and $\displaystyle \lim_{x \to -\infty} p(x) = \lim_{x \to \infty} p(x) = -\infty$. 
    \end{problem}

    \begin{solution}
      The ends of $p(x)$ are going the same direction, so the degree
      must be even. Since the limits are both $-\infty$, we also know that the leading coefficient
      should be negative. So one possible answer is 
      \fbox{$p(x)=-(x-2)^2(x-5)(x-7)$}.
    \end{solution}

  \item

    \begin{grading}
      3 points. There are many approaches, so this is a possible breakdown:
      \begin{itemize}[nosep]
      \item 1 for giving a correct possible formula for $p'$
      \item 1 for correctly going from their $p'$ to an $p$
      \item 1 for explanation
      \end{itemize}
      Note that there are infinitely many possible answers; any function of the form $a \left( \dfrac{x^4}{4} -2 x^2 \right) + b$ is correct. 
    \end{grading}

    \begin{problem}[\nstretch{1.25}]
      $p(x)$ is a degree 4 polynomial whose critical points are $x = 0$, $x = 2$, and $x = -2$. 
    \end{problem}

    \begin{solution}
      We know that the roots of $p'(x)$ are $x=0, \pm2$, so
      one possibility for $p'(x)$ is $p'(x)=x(x+2)(x-2)=x^3-4x$.
      A possible function with this derivative is 
      \fbox{$p(x)=\frac{1}{4}x^4 - 2x^2$.}
    \end{solution}

  \item

    \begin{grading}
      4 points:
      \begin{itemize}[nosep]
      \item 2 for giving a correct possible formula for $p''$: 0.5 for recognizing it's linear, 0.5 for having a root at $x = 1$, 1 for being positive on $(-\infty, 1)$ and negative on $(1, \infty)$
      \item 1 for correctly going from their $p''$ to an $p$
      \item 1 for explanation
      \end{itemize}
      Again, there are infinitely many possible answers; any function of the form $a \left( - \dfrac{x^3}{6} + \dfrac{x^2}{2} \right) + bx + c$ with $a > 0$ is correct. 
    \end{grading}

    \begin{problem}[\nstretch{1.25}]
      $p(x)$ is a cubic that's concave up on $(-\infty, 1)$ and concave down on $(1, \infty)$. 
    \end{problem}

    \begin{solution}
      There is an inflection point at $x=1$, and $p''(x)$ is positive
      for $x<1$ and negative for $>1$. So a possible second derivative
      is $p''(x)=-(x-1)=-x+1$. Based on this, a possible first derivative 
      is $p'(x)=-\frac{1}{2}x^2+x$. 
      Therefore a possible function 
      is \fbox{$p(x)=-\frac{1}{6}x^3+\frac{1}{2}x^2$}.
    \end{solution}
  \end{enumerate}

\item % 6.4.2

  \begin{grading}
    6 points:
    \begin{itemize}[nosep]
    \item 1 for goal
    \item 1 for labeling variables
    \item 0.5 for a correct formula for area in terms of their variables
    \item 1 for correctly writing an equation that says there are 600 feet of fencing
    \item 0.5 for using that equation to rewrite their area formula in terms of 1 variable
    \item 2 for finding where the maximum is and justifying that it really is a maximum. If they say the vertex is half-way between the roots but don't say anything that shows the vertex is a max rather than a min, please deduct 0.5. 
    \end{itemize}
  \end{grading}

  \begin{problem}[\vfill]    
    Troy is interested in skunks and has purchased two to study. He plans to fence off a rectangular pen for the skunks. He'll use an additional fence to divide the pen into 2 equal halves, one for each skunk. (See the picture for an example.) If Troy has 600 feet of fencing and wants to create the largest possible total area for his two skunks, what dimensions should he make the pen?

    \begin{tikzpicture}[x=1.5cm, y=1.5cm]
      \draw (0, 0) -- (3, 0) -- (3, 2) -- (0, 2) -- (0, 0); %
      \draw (1.5, 0) -- (1.5, 2); %
    \end{tikzpicture}
  \end{problem}

  \begin{solution}
  
    \textbf{Goal:} maximize area
    
    \textbf{Define variables:}
    Let $x =$ the width of the pen (meaning the divider is length $x$ as well) and $y =$ the length of the
    pen.
    
    \textbf{Formula for goal:} The total area is $A=xy$. To write this
    as a function of one variable, we can use information from the problem.
    Since Troy only has 600 feet of fencing, we know that $600 = 3x +2y$. Solving for $y$ gets us $y = 300-\frac{3}{2}x$. Therefore
    $A=x(300-\frac{3}{2}x)$.

    \textbf{Find maximum:}
    The area is a quadratic function in $x$. The coefficient of $x^2$
    would be negative if we expanded, the graph would be a concave down
    parabola. So there is a global maximum, at the vertex. The vertex is located half-way between the $x$-intercepts, 
    $0$ and $200$, so at $x=100$.

    \textbf{Answer the question:} When $x=100$, 
    $y=300-\frac{3}{2}(100)=150$. To maximize the area,
    Troy should make a pen with a width of $100$ feet and
    a length of $150$ feet.
  \end{solution}

\item  \begin{problem}[\newpage]
    The math department is interested in better understanding your math experience this semester. You can help by filling out \href{https://harvard.az1.qualtrics.com/jfe/form/SV_6PYrBw3v44khRvU}{this survey}. Your responses will be aggregated by a research team and anonymized before being shared with the teaching staff.
  \end{problem}

\end{enumerate}
\renewcommand{\psrnewpage}{{\centering\rule{5in}{0.4pt}}}

\psreflection{For next class, read \href{https://people.math.harvard.edu/~jjchen/mathma/docs/handouts/Connally-11-6.html}{\S 11.6 of \emph{Functions Modeling Change} by Connally et. al}.}
\end{document}
